\documentclass[11pt]{article}
\usepackage[a4paper,margin=1in]{geometry}
\usepackage{amsmath,amssymb,amsthm,microtype}
\usepackage{cleveref}
\newtheorem{theorem}{Theorem}[section]
\newtheorem{corollary}[theorem]{Corollary}
\newtheorem{proposition}[theorem]{Proposition}
\theoremstyle{remark}
\newtheorem{remark}[theorem]{Remark}
\newcommand{\R}{\mathbb R}
\newcommand{\TV}{\mathrm{TV}}
\title{A Sharp Physical-to-Affine Remainder for One Critical Endpoint Leg}
\author{Bin Wang}
\date{July 2026}
\begin{document}
\maketitle

\begin{abstract}
RH-323 propagated the critical endpoint profile through an exact affine
probability chain but left the physical curvature and row normalization
uncontrolled.  We now replace its first tangent leg by the exact
finite-noise folded kernel.  The physical row has an exact $L^1$ distance to
the full Gaussian centered at the curved quadratic image; this distance is
only a pair of state-boundary Gaussian tails and already includes the wrong
fold and exact row normalizer.  Removing the quadratic displacement then
costs order $\sigma$.  After averaging against the RH-322 entrance row we
obtain a fully explicit physical-to-affine joint bound and prove the sharp
phase-matched limit
\[
 \sigma^{-1}\|P_{\sigma,d}-A_d\|_1
 \longrightarrow
 \sqrt{2/\pi}\,u_c\{d^2+1+d\phi(d)/\Phi(d)\}>0.
\]
Thus the full affine remainder is not exponentially small.  Its local
$O(\sigma)$ scale is nevertheless smaller than the first-alias target scale.
No moving-order composition, second-leg remainder, parity cancellation, or
full trace replacement is asserted.
\end{abstract}

\section{The exact physical endpoint leg}

Let
\[
 f(x)=1-u_cx^2,
 \qquad r=u_c-1,
 \qquad \lambda=2u_cr,
 \qquad \alpha=2u_c,
\]
where
\[
 u_c=1.543689012692076\ldots,
 \quad r=0.543689012692076\ldots,
 \quad \lambda=1.678573510428322\ldots.
\]
The critical partition is $b=u_c^{-1/2}=0.8048595351\ldots$.
For $x,y\in[0,1]$, the archived folded row is
\begin{equation}\label{eq:physical-kernel}
 P_\sigma(x,y)=
 \frac{\phi_\sigma(y-f(x))+\phi_\sigma(y+f(x))}{Z_\sigma(x)},
 \qquad
 \phi_\sigma(t)=\sigma^{-1}\phi(t/\sigma),
\end{equation}
with $Z_\sigma(x)$ chosen so the row integrates to one.

Use the endpoint and repelling coordinates
\begin{equation}\label{eq:coordinates}
 x=1-\sigma v,
 \qquad y=r+\sigma u,
 \qquad
 I_\sigma=\left[-\frac r\sigma,\frac{1-r}{\sigma}\right].
\end{equation}
Direct substitution gives
\begin{equation}\label{eq:image-expansion}
 f(1-\sigma v)
 =-r+\sigma c_{\sigma,v},
 \qquad
 c_{\sigma,v}=\alpha v-u_c\sigma v^2.
\end{equation}
The exact rescaled physical row is therefore
\begin{equation}\label{eq:rescaled-row}
 K_\sigma(v,u)=
 \frac{
\phi(u+2r/\sigma-c_{\sigma,v})+\phi(u+c_{\sigma,v})
 }{Z_\sigma(1-\sigma v)}\mathbf1_{I_\sigma}(u),
\end{equation}
where, for every physical $0\le v\le\sigma^{-1}$,
\begin{equation}\label{eq:normalizer}
 Z_\sigma(1-\sigma v)
 =1-\overline\Phi\!\left(\frac{1-f(1-\sigma v)}\sigma\right)
   -\overline\Phi\!\left(\frac{1+f(1-\sigma v)}\sigma\right).
\end{equation}
Here $\phi$ and $\Phi$ are the standard normal density and distribution
function and $\overline\Phi=1-\Phi$.

On the endpoint branch $1-\sigma v\ge b$, the second lobe in
\eqref{eq:rescaled-row} is the physical folded image.  Define its full-line
curved Gaussian and the RH-323 tangent row by
\begin{equation}\label{eq:curved-tangent}
 C_{\sigma,v}(u)=\phi(u+c_{\sigma,v}),
 \qquad
 A_v(u)=\phi(u+\alpha v).
\end{equation}

\section{Exact row identities}

\begin{theorem}[Fold and normalization collapse to a boundary tail]
\label{thm:row-identities}
If $0\le v\le\sigma^{-1}$ and $1-\sigma v\ge b$, then
\begin{equation}\label{eq:exact-boundary}
 \boxed{
 \|K_\sigma(v,\cdot)-C_{\sigma,v}\|_1
 =2\left\{
 \overline\Phi\!\left(\frac r\sigma-c_{\sigma,v}\right)
 +\overline\Phi\!\left(\frac{1-r}\sigma+c_{\sigma,v}\right)
 \right\}.}
\end{equation}
Moreover,
\begin{equation}\label{eq:exact-shift}
 \|C_{\sigma,v}-A_v\|_1
 =4\Phi\!\left(\frac{u_c\sigma v^2}{2}\right)-2.
\end{equation}
Consequently the physical-to-tangent row distance is at most the sum of
\eqref{eq:exact-boundary} and \eqref{eq:exact-shift}.
\end{theorem}

\begin{proof}
The two unnormalized lobes in \eqref{eq:physical-kernel} have total mass one
on the positive half-line, so $0<Z_\sigma\le1$.  On $I_\sigma$, the numerator
of \eqref{eq:rescaled-row} contains $C_{\sigma,v}$ and division by
$Z_\sigma\le1$ only increases it.  Thus
$K_\sigma(v,u)\ge C_{\sigma,v}(u)$ on $I_\sigma$, while the physical row
vanishes outside.  The positive and negative $L^1$ parts are equal, and each
is the mass of $C_{\sigma,v}$ outside $I_\sigma$.  Its lower and upper tails
are precisely the two terms in \eqref{eq:exact-boundary}.

The centers of $C_{\sigma,v}$ and $A_v$ differ by $u_c\sigma v^2$.
Two unit Gaussians with center separation $t$ cross at their midpoint and
have $L^1$ distance $4\Phi(|t|/2)-2$, proving
\eqref{eq:exact-shift}.  The last assertion is the triangle inequality.
\end{proof}

This identity is stronger than separately bounding the reflected lobe and
$Z_\sigma-1$: their combined effect is exactly the discarded main-lobe
boundary mass.  The inequality
\begin{equation}\label{eq:linear-shift-bound}
 4\Phi(t/2)-2\le\sqrt{\frac2\pi}\,t,
 \qquad t\ge0,
\end{equation}
will isolate the curvature contribution.

\section{The finite-seed physical joint theorem}

Let the RH-322 entrance laws be
\begin{align}
 g_a(v)&=\frac{\phi(v-a)}{\Phi(a)}\mathbf1_{[0,\infty)}(v),
 \label{eq:limit-seed}\\
 \widetilde h_{\sigma,a}(v)
 &=\frac{\phi(v-a)+\phi(2L-v-a)}
 {\Phi(a)-\overline\Phi(2L-a)}\mathbf1_{[0,L]}(v),
 \qquad L=\sigma^{-1}.
 \label{eq:finite-seed}
\end{align}
Define the exact physical one-leg joint law and its affine target by
\begin{equation}\label{eq:joint-laws}
 P_{\sigma,a}(v,u)=\widetilde h_{\sigma,a}(v)K_\sigma(v,u),
 \qquad
 A_d(v,u)=g_d(v)A_v(u).
\end{equation}

Fix any
\begin{equation}\label{eq:eta-margin}
 0<\eta<1-b,
 \qquad
 m_\eta=u_c(1-\eta)^2-1>0.
\end{equation}
For reproduction we use the canonical midpoint
$\eta=(1-b)/2=0.0975702324\ldots$, for which
$m_\eta=0.2571486637\ldots$.

Put
\begin{equation}\label{eq:second-moment}
 M_2(a)=\int_0^\infty v^2g_a(v)\,dv
 =a^2+1+a\frac{\phi(a)}{\Phi(a)},
 \qquad
 c_0=\sqrt{\frac2\pi}\,u_c.
\end{equation}

\begin{theorem}[Explicit physical-to-affine joint remainder]
\label{thm:joint-bound}
For $\sigma>0$, $0\le a\le\sigma^{-1}$, and $d\ge0$,
\begin{align}
 \|P_{\sigma,a}-A_d\|_1
 \le{}& |a-d|
 +\frac{2\overline\Phi(\sigma^{-1}-a)}{\Phi(a)}
 +c_0\sigma M_2(a)
 \nonumber\\
 &+2\left\{
 \overline\Phi(m_\eta/\sigma)
 +\overline\Phi((1-r)/\sigma)
 \right\}
 +\frac{2\overline\Phi(\eta/\sigma-a)}{\Phi(a)}.
 \label{eq:explicit-joint-bound}
\end{align}
The same right side bounds the physical $U$-marginal error.  Uniformly for
$a,d$ in a fixed compact subset of $[0,\infty)$,
\begin{equation}\label{eq:compact-bound}
 \|P_{\sigma,a}-A_d\|_1
 \le |a-d|+C\sigma+Ce^{-c/\sigma^2}
\end{equation}
for sufficiently small $\sigma$, with $C,c>0$ depending only on the compact
set and $\eta$.
\end{theorem}

\begin{proof}
Extend $K_\sigma(v,\cdot)$ to all $v\ge0$ by setting it equal to $A_v$ for
$v>L$.  Markov lifting preserves the entrance $L^1$ distance, so RH-322 gives
\[
 \|\widetilde h_{\sigma,a}\widehat K_\sigma
      -g_a\widehat K_\sigma\|_1
 =\frac{2\overline\Phi(L-a)}{\Phi(a)}.
\]
The affine lift similarly gives
$\|g_aA-g_dA\|_1=\|g_a-g_d\|_1\le|a-d|$.

It remains to average the row error against $g_a$.  On
$0\le v\le\eta/\sigma$, the source stays in the negative endpoint branch.
The folded image has distance at least $m_\eta$ from zero, while its distance
from the upper state boundary is at least $1-r$.  Hence
\cref{thm:row-identities} and \eqref{eq:linear-shift-bound} give the first
three remainder terms in \eqref{eq:explicit-joint-bound} after using
$\int v^2g_a=M_2(a)$.  On the complement, two probability rows have distance
at most two, and
\[
 \Pr_{g_a}\{V>\eta/\sigma\}
 =\frac{\overline\Phi(\eta/\sigma-a)}{\Phi(a)}.
\]
This proves the explicit estimate.  Marginalization is an $L^1$ contraction.
Mills' inequality gives \eqref{eq:compact-bound}.
\end{proof}

\section{The curvature coefficient is sharp}

\begin{theorem}[Positive first-order obstruction]
\label{thm:sharp-coefficient}
For every fixed $d\ge0$,
\begin{equation}\label{eq:sharp-limit}
 \boxed{
 \lim_{\sigma\downarrow0}
 \frac{\|P_{\sigma,d}-A_d\|_1}{\sigma}
 =c_0M_2(d)
 =\sqrt{\frac2\pi}\,u_c
 \left\{d^2+1+d\frac{\phi(d)}{\Phi(d)}\right\}>0.}
\end{equation}
Thus the full physical-to-affine joint remainder is not exponentially small
and is not $o(\sigma)$.
\end{theorem}

\begin{proof}
Use the Markov extension from the preceding proof.  Replacing the exact
finite seed by $g_d$ changes the joint distance by at most
$2\overline\Phi(\sigma^{-1}-d)/\Phi(d)=o(\sigma)$.  Let
\[
 S_{\sigma,d}=\int_0^\infty g_d(v)
 \left\{4\Phi\!\left(\frac{u_c\sigma v^2}{2}\right)-2\right\}\,dv.
\]
On $v\le\eta/\sigma$, the reverse triangle inequality and
\eqref{eq:exact-boundary} show that the averaged physical row distance
differs from $S_{\sigma,d}$ by exponentially small boundary tails.  The
remaining entrance region also has exponentially small $g_d$ mass.

Pointwise,
\[
 \frac{4\Phi(u_c\sigma v^2/2)-2}{\sigma}
 \longrightarrow c_0v^2,
\]
and \eqref{eq:linear-shift-bound} dominates the quotient by $c_0v^2$.
Dominated convergence and \eqref{eq:second-moment} prove
\eqref{eq:sharp-limit}.
\end{proof}

The reproduction values for the dominant shift integral are
\[
\begin{array}{c|ccc|c}
d & \sigma=0.05 & 0.025 & 0.0125 & c_0M_2(d)\\ \hline
0   &1.227163&1.230543&1.231399&1.231686\\
0.5 &1.842222&1.850393&1.852473&1.853170
\end{array}
\]
where the three middle columns are $S_{\sigma,d}/\sigma$.  They reproduce
the analytic coefficient; they are not a fitted proof.

\section{First-alias scale and route firewall}

At the natural clock
\[
 k_\sigma=\frac{\log(1/\sigma)}{2\log\lambda}+O(1),
\]
the target radius $R=1.4$ gives
\begin{equation}\label{eq:alias-scale}
 R^{-2k_\sigma}=\Theta(\sigma^\theta),
 \qquad
 \theta=\frac{\log R}{\log\lambda}
 =0.6496301165\ldots<1.
\end{equation}
Consequently the phase-matched local remainder satisfies
\[
 \sigma=o(R^{-2k_\sigma}),
 \qquad
 k_\sigma\sigma=o(k_\sigma R^{-2k_\sigma}).
\]
This is a scalar scale comparison only.  It does not prove that a
moving-order product has bounded intermediate norms, that all legs satisfy
the same estimate, or that weighted trace contributions can be replaced by
separate absolute errors.

The implementation evaluates exact normalizers, row tails, Gaussian shift
distances, the explicit joint bound, and the sharp coefficient with
standard-library quadrature.  No exponent is fitted.  The result concerns a
forward probability density, not the peak-normalized backward Gaussian
observable and amplitude coefficient of RH-18.

\begin{remark}[Claim boundary]
The second physical leg, moving-order Duhamel composition, parity layer,
neighboring sibling/shell cancellation, branch-return matrix, and joint
first-alias trace law remain open.  No full-trace replacement or
$o(kR^{-2k})$ joint error is proved.  Gates A--E remain false/open.  This
paper constructs no Hilbert--Polya operator, identifies no Riemann zero,
proves no von Mangoldt trace formula or zeta-divisor equality, and does not
imply RH.
\end{remark}

RH-325 should state a moving-order Duhamel criterion that records retained
coordinates, phase transport, per-leg Markov contractions, and the precise
sum of local errors before any full-cycle application is attempted.

\end{document}
